%!TEX root = ../main.tex
\section{Closing the Loop: Industrial Data Recipes and Self-Evolving Engines}
\label{sec:recipes-engines}

This section covers the feedback-driven part of the operation axis and applies all four axes to composite systems.
The preceding sections examine one preparation operation at a time.
Industrial embodied models instead disclose a \emph{data recipe}: how several sources are transformed, assigned to training stages, and revised after evaluation.
These systems show which operations can be run at industrial scale, but several parts of a recipe usually change together.
A technical report may establish the performance of the complete model while leaving the mixture, preparation cost, or intermediate data versions undisclosed.
We therefore use a report as evidence about the disclosed system and do not assign its end-to-end gain to an unreported preparation step.

Figure~\ref{fig:industrial-recipe-map} aligns five public recipe families by source role, preparation mechanism, training use, and feedback.
We reconstruct these stages from primary reports, record only disclosed quantities, and identify conclusions that are our synthesis.
The comparison emphasizes functional structure instead of ranking reported performance numbers.
An hour of egocentric video, an hour of a 50-Hz robot trajectory, a simulated rollout, and a video--text clip are not equivalent data units.
Evaluations on different robots, tasks, model sizes, and budgets are not directly comparable either.

%!TEX root = ../../main.tex
\begin{figure*}[t]
\centering
\begin{tikzpicture}[
  font=\scriptsize,
  stage/.style={draw=black!40, rounded corners=2pt, text width=2.65cm,
    minimum height=2.9cm, align=left, inner sep=4pt},
  arr/.style={-{Latex[length=1.8mm]}, thick, draw=black!55},
  loop/.style={-{Latex[length=1.8mm]}, thick, dashed, draw=surveyrose}
]
\node[stage, fill=surveylightblue] (source) at (-6.3,0) {
  \textbf{1. Assign source roles}\\[2pt]
  robot control grounding\\
  human and web semantics\\
  simulation and generated coverage\\
  current-policy failures};
\node[stage, fill=surveylightteal, right=3.5mm of source] (align) {
  \textbf{2. Align representation}\\[2pt]
  schema and metadata\\
  reference frame and timing\\
  motion transfer\\
  inferred action uncertainty};
\node[stage, fill=surveylightorange, right=3.5mm of align] (check) {
  \textbf{3. Validate by role}\\[2pt]
  technical and cross-modal checks\\
  geometric or physical replay\\
  model and human audit};
\node[stage, fill=surveylightrose, right=3.5mm of check] (route) {
  \textbf{4. Route by stage}\\[2pt]
  broad pretraining\\
  embodiment alignment\\
  task specialization\\
  recovery post-training};
\node[stage, fill=surveylightgray, right=3.5mm of route] (evidence) {
  \textbf{5. Evaluate and revise}\\[2pt]
  component controls\\
  rollout outcomes\\
  retention and safety\\
  simulation and real tests};

\draw[arr] (source) -- (align);
\draw[arr] (align) -- (check);
\draw[arr] (check) -- (route);
\draw[arr] (route) -- (evidence);
\draw[loop] (evidence.south) -- ++(0,-0.55) -|
  node[below, pos=0.52, align=center, text=surveyrose, font=\footnotesize]
  {evaluation can revise source allocation, checks, and training roles}
  (source.south);
\end{tikzpicture}
\caption{Functional structure shared by public industrial data recipes. Systems differ in their sources and implementations, but meaningful comparison requires each stage to be reported separately.}
\label{fig:industrial-recipe-map}
\end{figure*}


\subsection{What the Industrial Recipes Actually Reveal}

The Qwen-RobotManip recipe shows why alignment must precede large-scale mixing.
Its corpus totals approximately 38,100 hours~\cite{yuan2026qwen_2606_17846}.
The disclosed total combines more than 11,000 hours of robot demonstrations, 1,933 hours of egocentric human video, and 24,808 hours of human-to-robot synthesis across 15 morphologies.
The scale is enabled by a chain that retargets human motion, reconstructs robot observations, and searches for feasible robot configurations.
It then applies five state--action checks and three cross-modal checks.
This is materially different from concatenating datasets and padding missing action dimensions with zeros.
A masked 80-dimensional schema fixes field meaning.
Camera-frame end-effector deltas align motion, while a positional encoding represents the camera view and prompts expose embodiment and control mode to the learner.
The report's audit found severe source-specific defects.
In one subset, its state--action trend test rejected 81\% of episodes.
This result is direct evidence that ``published robot data'' is not synonymous with ``training-ready robot data.''

The same conclusion appears in a different form in Gemini Robotics 1.5.
Merely adding other embodiments generally helps, but the reported single-embodiment, multi-embodiment-without-Motion-Transfer, and full-recipe ablation shows that a motion-alignment mechanism changes how much cross-robot data can be used ~\cite{team2025gemini_2510_03342}.
The benefit is asymmetric.
Data from other robots matters most for a low-data target, whereas a well-covered ALOHA platform has less to gain.
Cross-embodiment scale should report results for each source and target robot pair, not only a global hour count.

The $\pi_{0.5}$ recipe illustrates why training role can be more informative than source type.
It assigns robot data, high-level supervision, verbal instructions, and web data to different purposes.
Its broad pretraining mixture is mostly not the exact target household mobile-manipulation distribution.
Post-training then routes short successful episodes while retaining web and high-level supervision and adding verbal-instruction demonstrations~\cite{intelligence20250_2504_16054}.
A source can consequently be low value as an action target but high value as general semantic supervision or as supervision for high-level task structure.

$\pi_{0.7}$ pushes role separation further.
Instead of discarding every failure or suboptimal rollout, it records quality and strategy metadata and conditions the policy on the desired behavior~\cite{intelligence20260_2604_15483}.
In its ablations, autonomous evaluation data and metadata both matter.
Without contextual metadata, adding progressively lower-quality data can hurt, whereas the conditioned model continues to benefit in the reported setting.
This resolves an apparent conflict in the curation literature.  ``Keep only successes'' is reasonable when all retained actions are interpreted as expert targets.  ``Keep failures'' can be reasonable when outcome, mistake intervals, and desired behavior are represented and the objective can use them.
The disagreement is about \emph{routing and semantics}, not whether a failure bit has universal positive or negative value.

GR00T N1 illustrates a third pattern.
Synthetic expansion is useful only when later training continues to include verified real data and assigns generated data to a suitable role.
When a video lacks recorded controls, a latent-action model can encode its motion and an inverse-dynamics model can estimate the action between successive observations.
The recipe organizes web and human data, generated and simulated data, and real robot data as a pyramid.
It uses inferred latent or inverse-dynamics actions where recorded control is absent and real data for embodiment grounding~\cite{nvidia2025gr00t_2503_14734}.
In one generated-data post-training experiment it samples real and generated trajectories 1:1.
Its simulation branch keeps successful trajectories.
DreamGen follows four stages.
It adapts an image-to-video model with teleoperated robot video, generates task- and environment-conditioned video, recovers pseudo-actions with a latent-action or inverse-dynamics model, and post-trains the policy on the resulting video-derived trajectories ~\cite{jang2025dreamgen_2505_12705}.
Its DreamGenBench instruction and physics diagnostics correlate with downstream policy success in the reported experiments.
We interpret this result as support for routing at the dataset level, but it is not certification that every generated trajectory is executable.
The N1 report also supplies a valuable negative result.
Post-training without left-to-right handover examples removed that capability.
Retaining broad real data and setting minimum sampling shares for important sources can reduce forgetting during later training stages.

Qwen-RobotWorld expands the generator side with 8.6 million clips and more than 200 million frames in a corpus spanning more than 20 embodiments and 500 action categories ~\cite{zhang2026qwen_2606_17030}.
Its processing chain first standardizes video, then produces five-layer action descriptions and sends low-scoring captions from an LLM-based judge to targeted prompt revision and human review.
Training preserves a 70:30 mixture of embodied and general data and progressively introduces single-view, multi-view, concatenated-view, and complex cross-domain data.
The report also excludes AI-generated media from the retained general-data mixture because of artifact and bias concerns.
This choice is a useful counterexample to indiscriminate synthetic scaling.
Language can describe human, robot, and simulated motion through one interface, but it does not recover calibrated joint actions or certify contact.
Outputs should be routed by training role.
Broad video dynamics, proposal, and planning data need different validation from behavior-cloning targets.
A world model may support virtual evaluation, but the same checkpoint should not generate and uncritically approve its own samples.

The Gemini Robotics reports also show that disclosure is part of what can be evaluated about a recipe.
They describe thousands of hours of ALOHA~2 demonstrations, multimodal data without action labels, and 2,000--5,000 high-quality episodes per specialization task ~\cite{team2025gemini_2503_20020}.
Version 1.5 reports thousands of tasks over three embodiments and interleaved real-robot comparisons.
Its Motion Transfer ablation isolates the alignment mechanism, while its exact source quantities, filters, and annotation lineage remain closed.
The same report notes that more than 90\% of development evaluation episodes ran in simulation and that final model quality still requires real testing~\cite{team2025gemini_2510_03342}.
This is good systems practice, but it also means simulated iteration speed and real deployment evidence must remain separate fields.
Reproducibility should be reported separately for data, preparation code, training mixture, model, and evaluation.
A single ``open'' or ``closed'' label hides these differences.

Table~\ref{tab:industrial-recipes} condenses the preceding analysis into a reference index.
Its evidence column records what can be concluded from the public report, not a score for the complete recipe.

\begin{table*}[t]
\centering
\caption{Structural comparison of public industrial data recipes. Quantities use each report's own units and are not comparable across rows.}
\label{tab:industrial-recipes}
\scriptsize
\setlength{\tabcolsep}{4pt}
\begin{tabularx}{\textwidth}{@{}p{3.0cm}p{4.0cm}p{5.1cm}X@{}}
\toprule
System & Source roles & Main preparation decision & Evidence boundary \\
\midrule
$\pi$ family~\cite{intelligence20250_2504_16054,intelligence20260_2604_15483,intelligence20250_2511_14759}
& Cross-robot, web, mixed-quality, and target-task data serve different training stages.
& Normalize actions, route sources by stage, and condition on quality or outcome metadata.
& Family ablations support several choices, but complete volumes and mixture weights are not disclosed. \\
GR00T family~\cite{nvidia2025gr00t_2503_14734,jang2025dreamgen_2505_12705,nvidia2026gr00t17}
& Web and human video broaden semantics; simulation and generation expand trajectories; real data ground control.
& Infer missing actions, retain validated trajectories, and mix generated data with real demonstrations.
& Within-report studies show utility and forgetting; product-scale figures are engineering evidence. \\
Gemini Robotics~\cite{team2025gemini_2503_20020,team2025gemini_2510_03342}
& Multimodal pretraining, three robot embodiments, and narrow specialization sets have distinct roles.
& Relabel for embodied reasoning and use Motion Transfer to learn from heterogeneous robot motion.
& The report ablates embodiment diversity and Motion Transfer; source quantities and filtering remain largely closed. \\
Qwen-RobotManip~\cite{yuan2026qwen_2606_17846}
& Robot demonstrations, human video, and human-to-robot synthesis form a roughly 38,100-hour corpus.
& Retarget human motion, apply state--action and cross-modal checks, and align representation, motion, and behavior.
& The report ablates major components, but does not isolate the physical error rate of all synthetic data. \\
Qwen-RobotWorld~\cite{zhang2026qwen_2606_17030}
& Embodied video supplies dynamics; general multimodal data preserves broad knowledge.
& Standardize video, annotate actions in language, review low-scoring labels, and use a staged curriculum.
& Generation and prediction are demonstrated; per-clip action validity and policy utility are not guaranteed. \\
\bottomrule
\end{tabularx}
\end{table*}


\subsection{When a Pipeline Becomes a Data Engine}

A static pipeline applies a fixed sequence of operations to a raw dataset.
A data engine differs because the performance of the current learner changes which data are collected, generated, labeled, weighted, or checked in a later round.
In each round, the system trains or updates a model on a versioned training distribution and evaluates that model on a declared target.
The evaluation records task outcomes together with interventions, uncertainty, and coverage gaps.
A controller uses those observations to choose the next predefined data action within explicit human, compute, simulation, and robot budgets.

The feedback requirement matters.
Re-running the same loading and cleaning script is workflow automation, and periodically fine-tuning on every new log is a model-update schedule.
Neither is a data engine unless results from the current learner determine which records are sought, changed, retained, or assigned to a new training role.

Table~\ref{tab:data-engines} compares representative loops through five controller fields.
This exposes differences hidden by the common ``flywheel'' label.
RoboCat and AutoRT close acquisition loops through practice and foundation-model task proposal~\cite{bousmalis2023robocat_2306_11706,ahn2024autort_2401_12963}.
DexFlyWheel is a simulation-centered generative loop seeded by one human demonstration~\cite{zhu2025dexflywheel_2509_23829}.
The $\pi^{*}_{0.6}$ report calls its experience-learning loop RECAP.
It combines outcome-derived value estimates with autonomous trials and optional human corrections.
The Scalable Online Post-Training system (SOP) and Learning While Deploying (LWD) are also post-training loops, but their data semantics and update cadence differ.
Robot-Powered Data Flywheels uses deployment to adapt an upstream perception model, showing that the updated learner need not be the action policy ~\cite{grannen2025robot_2511_19647}.
The term names a systems pattern, not one algorithm.

\begin{table*}[t]
\centering
\caption{Data engines compared as feedback controllers.  The discriminating
questions are what detects a gap, which data action follows, how the learner is
updated, and what prevents a locally improving loop from drifting or collapsing.}
\label{tab:data-engines}
\scriptsize
\setlength{\tabcolsep}{3.5pt}
\begin{tabularx}{\textwidth}{@{}p{2.8cm}p{3.5cm}p{6.15cm}X@{}}
\toprule
System & Gap/feedback signal & Data action and model update & Control mechanism and remaining limit \\
\midrule
RoboCat / AutoRT
~\cite{bousmalis2023robocat_2306_11706,ahn2024autort_2401_12963}
& New-task practice outcomes; foundation-model task proposals, affordances,
and safety checks.
& Generate practice or fleet experience, add it to a generalist mixture, and
fine-tune; AutoRT separates semantic task proposal from executable robot
collection.
& Safety rules and human/environment constraints gate collection.  The loop
does not autonomously attribute gains among proposer, data, objective, and
model changes. \\
\midrule
$\pi^{*}_{0.6}$/RECAP
~\cite{intelligence20250_2511_14759}
& Terminal success/failure labels, time-to-outcome values, autonomous trials,
and optional human interventions.
& Train a language-conditioned value function on accumulated experience and
extract an advantage-conditioned VLA; repeat collection and offline updates.
& Policy/value are restarted from the audited pretrained checkpoint at each
iteration in the reported implementation to limit drift.  Rewards and
interventions still require task-specific human infrastructure. \\
\midrule
DexFlyWheel
~\cite{zhu2025dexflywheel_2509_23829}
& Simulator reward and rollout success under new objects/configurations.
& One human seed bootstraps imitation learning; residual RL explores; successful
rollouts and geometric augmentation form the next iteration's dataset.
& Closed-loop simulation makes validity testable, but manual reward design and
the digital twin define the reachable distribution; it is not an open-world
failure-discovery system. \\
\midrule
SOP
~\cite{pan2026sop_2601_03044}
& Fresh on-policy failures/interventions and sliding-window online versus
offline training losses.
& A distributed actor fleet streams complete episodes to one learner; a
task-balanced sampler mixes static and online buffers and plug-in HG-DAgger or
RECAP updates are broadcast asynchronously.
& Uniform task weights and a clipped 0.2--0.8 online ratio preserve coverage;
checkpoints change only between episodes.  It still depends on human
interventions or task rewards and is tested over hours, not lifelong operation. \\
\midrule
LWD
~\cite{wang2026learning_2605_00416}
& Sparse terminal rewards over successes, failures, retries, partial progress,
and intervention segments from a 16-robot fleet.
& Distributional value learning uses heterogeneous offline/online replay;
critic gradients update a flow-based action expert through adjoint matching,
and the shared policy is repeatedly redeployed.
& A fixed offline reference, clipped double critic, mixed replay, and freezing
the VLM backbone online reduce instability.  Four-hour studies do not yet
establish long-term verifier or distribution stability. \\
\midrule
Robot-Powered Data Flywheels
~\cite{grannen2025robot_2511_19647}
& Deployment novelty and errors, with environmental/catalog signals that can
provide labels during use.
& A two-week library deployment collects task-relevant robot observations and
uses them to adapt an upstream visual foundation model, feeding improved
perception back to the robot system.
& Demonstrates that a flywheel may update perception rather than only the
action policy.  Site-specific labels and limited duration leave transfer and
long-run bias unresolved. \\
\midrule
DEED
~\cite{sis2026closing_2607_20345}
& Deployment experience and failures reveal gaps in an offline demonstration
set.
& Experience is converted into new or revised demonstrations, then used to
update the policy and the next data-collection decision.
& Makes demonstration evolution explicit, but causal promotion tests,
independent verifiers, and rollback remain necessary when several loop
components change together. \\
\bottomrule
\end{tabularx}
\end{table*}


\subsection{Reinforcement Learning, Post-Training, and Failure-Driven Evolution}

Post-training changes preparation because the deployed policy determines which states are observed.
Offline demonstration data approximate an expert or collector distribution.
Data collected by the current policy, commonly called on-policy data, reveal the states and errors that model itself produces.
Their value cannot be compared by hours alone.
SOP reports that, for one base model and task family, three hours of targeted on-policy interaction closed substantially more of the gap than an additional 80 hours of static demonstrations~\cite{pan2026sop_2601_03044}.
This is a within-paper finding, not a universal exchange rate.
It supports the narrower causal claim that failure-targeted experience can have high marginal value near deployment.

The key operation is not ``add failures,'' but to convert a rollout into training units with explicit roles.
A successful segment may be an imitation target.
A valid prefix of a failed episode can train reaching or value prediction.
An intervention provides both a takeover boundary and a corrective action.
A near miss can be a negative example for a model that scores task outcomes, while an unsafe action can calibrate a safety verifier but must not silently enter behavior cloning.
Ambiguous cases should be retained with uncertainty and routed to review.
RoboReward makes this distinction concrete.
It changes outcome labels and clips trajectories in time to construct negative and near-miss examples for a specialized reward VLM ~\cite{lee2026roboreward_2601_00675}.
Its real-robot study also cautions that a general embodied-reasoning model is not automatically a calibrated reward model for the target data distribution.

RECAP combines demonstrations, autonomous rollouts, terminal outcomes, and human corrections.
Its value model predicts time to success or a large failure penalty.
The policy then favors actions estimated to improve over the reference behavior ~\cite{intelligence20250_2511_14759}.
Notably, the reported multi-round implementation fine-tunes each value model and policy from the common pretrained checkpoint instead of the previous iteration, which the authors found useful for limiting drift.
SOP reduces the latency between failure and correction with distributed actors, balances tasks uniformly, clips each task's online fraction to $[0.2,0.8]$, and updates robot checkpoints only between episodes ~\cite{pan2026sop_2601_03044}.
LWD instead models a distribution of possible returns rather than a single expected return and replays both offline and online experience.
During online updates, it freezes the VLM backbone and changes only the action-producing module, another explicit stability choice ~\cite{wang2026learning_2605_00416}.
These mechanisms show that data and model updates must be designed together.
Buffer composition, reward semantics, update latency, and trainable modules all determine whether a failure is useful.

Three evaluation controls are still rare.
A multi-round system should report utility throughout the sequence of robot interactions, not only at its final checkpoint.
It should audit retained broad and safety capabilities because correcting today's failure can erase yesterday's rare skill.
It should also compare the proposed new data with a volume- and compute-matched sample from the same collection round before promoting a successor model.
Without that control, a better successor may reflect more interaction, more updates, or a changed reward rather than the engine's data decision.

\subsection{From Static Curation to a Data-Preparation Agent}

Figure~\ref{fig:autonomy-ladder} organizes the evolution by \emph{decision authority}.
At L0, people choose a fixed pipeline.
L1 optimizes an offline subset or mixture for a learner.
L2 uses rollout failures or interventions for one or a few post-training rounds.
An L3 engine repeatedly executes a predesigned collect--verify--route--train--deploy controller.
An L4 data-preparation agent must additionally choose which operator or experiment to run, estimate its information value and cost, decide when evidence is sufficient, and stop or roll back under risk.
Existing systems demonstrate substantial pieces through L3.
Calling them fully autonomous L4 agents would overstate the evidence.

\begin{figure*}[t]
\centering
\begin{tikzpicture}[
  level/.style={draw, rounded corners=2pt, align=center, text width=2.88cm,
    minimum height=1.75cm, inner sep=4pt, font=\footnotesize},
  arr/.style={-{Latex[length=1.8mm]}, thick, draw=black!55},
  axis/.style={-{Latex[length=2mm]}, very thick, draw=surveyblue}
]
\node[level, fill=surveylightgray] (l0) {\textbf{L0: Static}\\fixed rules and one-shot mixture\\\emph{people choose the pipeline}};
\node[level, fill=surveylightblue, right=3mm of l0] (l1) {\textbf{L1: Learner-aware}\\loss, influence, or proxy selection\\\emph{offline adaptation}};
\node[level, fill=surveylightteal, right=3mm of l1] (l2) {\textbf{L2: Failure-driven}\\rollouts alter later training data\\\emph{one or a few rounds}};
\node[level, fill=surveylightorange, right=3mm of l2] (l3) {\textbf{L3: Data engine}\\repeated collect--verify--train loop\\\emph{fixed controller}};
\node[level, fill=surveylightrose, right=3mm of l3] (l4) {\textbf{L4: Preparation agent}\\chooses operators, evidence, and stopping\\\emph{adaptive controller}};
\draw[arr] (l0) -- (l1);
\draw[arr] (l1) -- (l2);
\draw[arr] (l2) -- (l3);
\draw[arr] (l3) -- (l4);
\draw[axis] ([yshift=-7mm]l0.south west) -- node[below, align=center, font=\small] {increasing decision authority and target relevance\\increasing feedback dependence, evidence cost, and safety risk} ([yshift=-7mm]l4.south east);
\end{tikzpicture}
\caption{Evolution of data preparation by decision authority. Existing systems demonstrate substantial components through L3. L4 additionally requires autonomous experiment choice, matched deployment tests, stopping, and rollback.}
\label{fig:autonomy-ladder}
\end{figure*}


An L4 agent differs from a prescribed loop because it chooses the next experiment.
At each round, it needs versioned data and model lineage, calibrated evaluation, a budget, and explicit safety limits.
It may then choose whether to obtain new records, revise existing units, change their training role, run a diagnostic, or stop.
These choices produce evidence at different delays and costs.
A technical check may return immediately, while controlled policy training or a robot rollout may take hours and expose physical risk.
The agent should optimize downstream policy performance over the full sequence of decisions after accounting for preparation cost and risk.
It should not maximize an immediate score produced by a judge that changes across rounds.

Moving from L3 to L4 requires more than automation.
Every operator needs an explicit input and output schema and a reversible record of its effects.
The agent must represent uncertainty about policy improvements that are observed only after training and use controlled experiments that isolate one material data decision at a time.
Governance must allow abstention, preserve fixed reference data and capabilities, run small controlled training checks, and roll back data and model versions together.
An LLM can serve as a planner or semantic interface, but language fluency does not supply physical validation or permission for risky exploration.
The step from L3 to L4 is the construction of an auditable experimental controller whose authority expands only when its evidence and safety controls support it.
